\centerline{\bf \Huge Оптимизация. КР}
\centerline{\bf \Huge }


\problem{\textbf{1}} Предприниматель купил здание и собирается открыть в нём отель. В отеле могут быть
стандартные номера площадью 27 квадратных метров и номера «люкс» площадью 45
квадратных метров. Общая площадь, которую можно отвести под номера, составляет 981
квадратный метр. Предприниматель может поделить эту площадь между номерами различных
типов, как хочет. Обычный номер будет приносить отелю 2000 рублей в сутки, а номер «люкс» —
4000 рублей в сутки. Какую наибольшую сумму денег сможет заработать в сутки на своём отеле
предприниматель?

\problem{\textbf{2}} В распоряжении прораба имеется бригада рабочих в составе 26 человек. Их нужно распределить
на строительство двух частных домов, находящихся в разных городах.
Если на строительстве первого дома $t$ работает человек, то их суточная зарплата составляет $3t^2$
денежных единиц. Если на строительстве второго дома работает $t$ человек, то их суточная
зарплата составляет $4t^2$ денежных единиц. Дополнительные суточные накладные расходы, то
есть транспорт, питание и тому подобное, обходятся в 4 денежных единицы в расчёте на одного
рабочего при строительстве первого дома и в 3 денежных единицы при строительстве второго
дома.
Как нужно распределить на эти объекты рабочих бригады, чтобы все выплаты на их суточное
содержание, то есть суточная зарплата и суточные накладные расходы, оказались
наименьшими? Сколько денежных единиц в сумме при таком распределении составят все
суточные затраты, то есть зарплата и накладные расходы?

\problem{\textbf{3}} Строительство нового завода стоит 159 млн рублей. Затраты на производство $x$ тыс. ед.
продукции на таком заводе равны $0,5x^2 + 2x + 6$ млн рублей в год. Если продукцию завода
продать по цене $p$ тыс. рублей за единицу, то прибыль фирмы в млн рублей за один год составит
$px - (0,5x^2 + 2x + 6)$. Когда завод будет построен, фирма будет выпускать продукцию в таком
количестве, чтобы прибыль была наибольшей. При этом в первый год $p$ = 10, а далее каждый
год возрастает на 1. За сколько лет окупится строительство?

.

.

.


\problem{\textbf{4}} В двух шахтах добывают алюминий и никель. В первой шахте имеется 60 рабочих, каждый
из которых готов трудиться 5 часов в день. При этом один рабочий за час добывает 2 кг
алюминия или 3 кг никеля. Во второй шахте имеется 260 рабочих, каждый из которых готов
трудиться 5 часов в день. При этом один рабочий за час добывает 3 кг алюминия или 2 кг
никеля.
Обе шахты поставляют добытый металл на завод, где для нужд промышленности производится сплав алюминия и никеля, в котором на 2 кг алюминия приходится 1 кг никеля. При этом
шахты договариваются между собой вести добычу металлов так, чтобы завод мог произвести
наибольшее количество сплава. Сколько килограммов сплава при таких условиях ежедневно
сможет произвести завод?